\chapter{Control: Managing the Error State}

\otbookepigraph{Control is just the art of making the universe behave like your math model.}{OT Interpretation}

\section{De-construction: PD Control}
The standard PD law is:
\begin{equation}
    \tau = \sym{kp} \sym{error-e} + \sym{kd} \dot{\sym{error-e}}
\end{equation}

\section{Re-interpretation: The ``Spring-Damper'' Analogy}
\begin{intuitionbox}[Virtual Physics]
    Don't see \sym{kp} and \sym{kd} as gain numbers. See them as virtual components.
    \begin{itemize}
        \item \sym{kp} is a \textbf{Virtual Spring} that pulls the robot back to the target.
        \item \sym{kd} is a \textbf{Virtual Shock Absorber} that stops it from oscillating wildly.
    \end{itemize}
    Tuning a robot is just choosing how stiff and how squishy you want your virtual world to be.
\end{intuitionbox}

\section{Computed Torque Control}
\begin{designrule}[The Cancellation Move]
    The most powerful move in robotics is \textbf{Inverse Dynamics}. We use our known model (\sym{mass-matrix}, \sym{coriolis}, \sym{gravity-g}) to calculate and subtract the nonlinearities. This leaves behind a perfectly simple linear system.
\end{designrule}

\section{Mastery Check}
\begin{practice}[The Stability Question]
    Why do we often add a gravity compensation term \(\hat{g}(q)\) to our PD control? What happens to the \sym{error-e} in steady-state if we don't? (It sags due to gravity).
\end{practice}
